\section{Conclusion}


In this work, we introduced reference-adjusted evaluation, a new paradigm for LLM-as-judge summarization evaluation that bridges the gap between the expensive reference-based setting and the weaker reference-free setting. Our experiments across three open-source LLMs showed that reference-adjusted evaluation improves over the reference-free setting and, in some cases, approaches or outperforms the gold reference-based setting, especially when only a single gold reference is available. These results suggest that researchers can meaningfully improve their evaluations on new test sets without incurring the full cost of human annotation by leveraging either machine generated summaries or naturally occurring data. Our reference-adjusted paradigm provides a practical and scalable path toward more robust automatic evaluation of summarization systems.